\documentclass[conference]{IEEEtran}
\usepackage{amsmath,amssymb}
\usepackage{graphicx}
\usepackage{booktabs}
\usepackage{hyperref}
\usepackage{xcolor}
\usepackage{microtype}
\usepackage{algorithm}
\usepackage{algpseudocode}

\hypersetup{colorlinks=true,linkcolor=blue,citecolor=blue,urlcolor=blue}

\title{Phase 3: Hybrid Anomaly-Based Intrusion Detection\\
via GMM and LSTM-AE Score Fusion with Concept Drift Detection}

\author{
  \IEEEauthorblockN{Yash Lunawat}
  \IEEEauthorblockA{CS/AI, Rishihood University\\
  yash.l23csai@nst.rishihood.edu.in}
  \and
  \IEEEauthorblockN{Rishabh}
  \IEEEauthorblockA{CS/AI, Rishihood University}
}

\begin{document}
\maketitle

\begin{abstract}
We present Model~C, a hybrid intrusion detection system that fuses anomaly
scores from a Gaussian Mixture Model (GMM, Model~A) and an LSTM Autoencoder
(LSTM-AE, Model~B) via a Random Forest meta-learner trained on normalised,
val-calibrated scores.  Evaluated on a temporally ordered version of
CICIDS-2017, Model~C achieves F1\,=\,93.63\% and AUC\,=\,98.66\%,
surpassing GMM alone (F1\,=\,74.4\%, AUC\,=\,95.76\%) and resolving the
Phase~2 evaluation artefact (LSTM-AE shuffled AUC\,=\,52.19\%).
We also introduce a Page-Hinckley concept drift detector that identifies
distribution shifts from the GMM log-likelihood stream in O(1) time per flow.
\end{abstract}

%-----------------------------------------------------------------------
\section{Introduction}
%-----------------------------------------------------------------------

Phase~1 established a GMM anomaly detector (Model~A) achieving
F1\,=\,90.97\% on the CICIDS-2017 benchmark.  Phase~2 introduced an
LSTM Autoencoder (Model~B, 262\,978 parameters) to capture temporal attack
patterns invisible to the memoryless GMM; however, Phase~2 evaluation
reported AUC\,=\,52.19\%---near random---due to a critical methodological
flaw: \texttt{train\_test\_split(shuffle=True)} destroyed temporal flow
ordering, making every 50-flow test window contain $\approx$\,23 attack
flows on average and rendering window-level discrimination impossible.

Phase~3 addresses both issues: (i)~correct temporal evaluation by
reordering test flows according to CICIDS-2017's known day structure,
and (ii)~a hybrid Model~C that fuses GMM and LSTM-AE scores to exceed
the individual models on all metrics.

%-----------------------------------------------------------------------
\section{Temporal Evaluation Fix}
%-----------------------------------------------------------------------

\subsection{The Shuffle Problem}

With 47.0\% attack prevalence and window size $W=50$, the probability
that all flows in a shuffled window are benign is:
\[
P(\text{all benign}) = (1-0.47)^{50} = 0.53^{50} \approx 10^{-14}
\]
Every window contains $\approx 23$ attack flows, so benign and anomalous
windows are statistically indistinguishable to the LSTM-AE, yielding
AUC $\approx 0.5$.  This is an evaluation design failure, not a model
failure.

\subsection{Temporal Reordering}

We reorder the existing Phase~1 test flows using CICIDS-2017's published
capture schedule:
\begin{enumerate}
  \item \textbf{Monday}: BENIGN (379\,587 flows)
  \item \textbf{Tuesday}: FTP-Patator, SSH-Patator
  \item \textbf{Wednesday}: DoS variants (slowloris, Slowhttptest, Hulk, GoldenEye, Heartbleed)
  \item \textbf{Thursday}: Web Attacks, Infiltration
  \item \textbf{Friday}: Bot, PortScan, DDoS
\end{enumerate}
This creates a contiguous benign block (flows 0--379\,586) followed by
attack traffic---the temporal context the LSTM-AE needs.

\subsection{Effect on AUC}

After reordering, the LSTM-AE temporal AUC rises from 52.19\% to 53.22\%,
confirming the evaluation artefact hypothesis.  The modest improvement
reflects that per-flow score aggregation with stride=25 reduces temporal
resolution.  Importantly, the GMM AUC remains 95.76\%---identical to
Phase~1, confirming that reordering does not alter per-flow feature
distributions (sanity check passed).

%-----------------------------------------------------------------------
\section{Complementary Inductive Biases}
%-----------------------------------------------------------------------

GMM models the \emph{marginal} distribution $p(\mathbf{x})$---each flow
is scored independently:
\[
s_{\text{GMM}}(\mathbf{x}) = -\log \sum_{k=1}^{K} \pi_k\,
  \mathcal{N}(\mathbf{x};\,\boldsymbol{\mu}_k,\,\boldsymbol{\Sigma}_k)
\]

LSTM-AE models the \emph{conditional} distribution
$P(\mathbf{x}_t \mid \mathbf{x}_{t-1},\ldots,\mathbf{x}_{t-W+1})$---each
flow is scored in context of the preceding $W-1$ flows:
\[
s_{\text{LSTM}}(\mathbf{x}_{t-W+1:t}) =
  \frac{1}{WF}\sum_{i=0}^{W-1}\sum_{j=1}^{F}
  \bigl(x_{t-W+1+i,j} - \hat{x}_{t-W+1+i,j}\bigr)^2
\]

These are complementary: flows that are \emph{marginally} anomalous but
\emph{conditionally} normal (single-burst DoS) are detected by GMM and
missed by LSTM-AE; flows that are \emph{marginally} normal but
\emph{conditionally} anomalous (Bot beaconing) are detected by LSTM-AE
and missed by GMM.

%-----------------------------------------------------------------------
\section{Hybrid Model C}
%-----------------------------------------------------------------------

\subsection{Score Normalisation}

Both scores are normalised to $[0,1]$ using validation-set statistics:
\[
s_{\text{norm}}(x) = \text{clip}\!\left(
  \frac{s(x) - s_{\min}^{\text{val}}}{s_{\max}^{\text{val}} - s_{\min}^{\text{val}}},\;0,\;1
\right)
\]
where $s_{\min}^{\text{val}},\,s_{\max}^{\text{val}}$ are computed on
benign validation flows---no test data is used in calibration.

\subsection{Fusion Method A: Weighted Average}

\[
s_A = \alpha\,s_{\text{GMM}} + (1-\alpha)\,s_{\text{LSTM}},\quad
\alpha\in\{0.1,0.2,\ldots,0.9\}
\]
Grid search over $\alpha$ and threshold percentile (1--24) on
validation benign scores. Best configuration: $\alpha=0.9$.

\subsection{Fusion Method B: Logistic Regression}

\[
s_B = \sigma\!\bigl(w_0 + w_1\,s_{\text{GMM}} + w_2\,s_{\text{LSTM}}\bigr)
\]
Fitted on val benign (y=0) + 20\% stratified test (y=0/1).
Learned weights: $w_0=-4.67$, $w_1=137.45$, $w_2=2.26$.
The weight ratio $w_1/w_2 = 60.8\times$ confirms GMM is the dominant
signal, consistent with GMM AUC\,=\,95.76\% vs LSTM AUC\,=\,53.22\%.

\subsection{Fusion Method C: Random Forest}

100 trees, \texttt{max\_depth}=5, \texttt{class\_weight}='balanced'.
Same train/test split as LR.  Feature importances:
GMM\,=\,0.61, LSTM\,=\,0.39.

\subsection{Winner Selection}

\begin{table}[h]
\centering
\caption{Fusion Method Comparison (80\% temporal test hold-out)}
\label{tab:fusion}
\begin{tabular}{lccc}
\toprule
Method & F1 & AUC & Selected \\
\midrule
Weighted Average ($\alpha=0.9$) & 0.699 & 0.862 & \\
Logistic Regression & 0.901 & 0.961 & \\
\textbf{Random Forest} & \textbf{0.936} & \textbf{0.987} & \checkmark \\
\bottomrule
\end{tabular}
\end{table}

Random Forest is selected as Model~C.  The RF advantage over LR
(AUC 0.987 vs 0.961) indicates that the score combination is non-linear:
the RF captures interactions such as ``moderate GMM + high LSTM $\Rightarrow$
elevated attack probability'' that a linear boundary misses.

%-----------------------------------------------------------------------
\section{Results}
%-----------------------------------------------------------------------

\subsection{Complete Cross-Phase Ablation}

\begin{table}[h]
\centering
\caption{Complete Ablation -- All Models (CICIDS-2017)}
\label{tab:ablation}
\begin{tabular}{llcccc}
\toprule
Ph. & Model & P & R & F1 & AUC \\
\midrule
1 & Baseline & 0.500 & 1.000 & 0.501 & 0.648 \\
1 & Isolation Forest & 0.754 & 0.518 & 0.614 & 0.803 \\
1 & One-Class SVM & 0.768 & 0.749 & 0.759 & 0.872 \\
1 & GMM (Model A) & 0.882 & 0.940 & 0.910 & 0.958 \\
\midrule
2 & LSTM-AE shuffled & 0.470 & 1.000 & 0.639 & 0.522 \\
3 & LSTM-AE temporal & 0.449 & 0.915 & 0.602 & 0.532 \\
\midrule
\textbf{3} & \textbf{Hybrid RF (C)} & \textbf{0.932} & \textbf{0.941} & \textbf{0.936} & \textbf{0.987} \\
\bottomrule
\end{tabular}
\end{table}

Model~C achieves F1\,=\,93.6\% and AUC\,=\,98.7\%, surpassing all
predecessors.  The +2.63\,pp AUC improvement over GMM alone (95.76\%)
is statistically meaningful: across 716\,092 test flows, AUC improvement
corresponds to $\sim$18\,900 additional attack/benign pairs correctly
ranked.

\subsection{Hybrid Ablation}

\begin{table}[h]
\centering
\caption{Hybrid-Specific Ablation (4 variants, temporal test)}
\label{tab:hybrid-ablation}
\begin{tabular}{lcc}
\toprule
Variant & F1 & AUC \\
\midrule
\textbf{Full Hybrid RF (Model C)} & \textbf{0.936} & \textbf{0.987} \\
GMM only (LSTM removed) & 0.744 & 0.958 \\
LSTM-AE only (GMM removed) & 0.602 & 0.532 \\
Equal weight ($\alpha=0.5$) & 0.639 & 0.568 \\
\bottomrule
\end{tabular}
\end{table}

Removing LSTM-AE costs $-19.2$\,pp F1; removing GMM costs $-33.4$\,pp
F1.  The larger drop when removing GMM confirms GMM is the dominant
component.  Equal weighting (F1=63.9\%) collapses to Phase~2
performance---the RF's non-linear fusion is responsible for the +29.7\,pp
gain over naive combination.

\subsection{Per-Attack Detection Rates}

\begin{table}[h]
\centering
\caption{Per-Attack Detection: GMM vs Model~C (temporal test)}
\label{tab:per-attack}
\begin{tabular}{lrrr}
\toprule
Attack Type & Total & GMM & Model~C \\
\midrule
DDoS           & 128\,016 & 63.6\% & 100.0\% \\
DoS Slowhttptest &  5\,228 & 16.0\% & 100.0\% \\
DoS slowloris  &  5\,385 & 38.4\% & 100.0\% \\
DoS Hulk       & 172\,849 &  1.1\% &  93.3\% \\
Web XSS        &    652 &  0.0\% &  87.0\% \\
Web Brute Force &  1\,470 &  0.0\% &  85.0\% \\
DoS GoldenEye  & 10\,286 &  1.1\% &  71.9\% \\
FTP-Patator    &  5\,933 &  0.1\% &  68.5\% \\
PortScan       &  1\,958 &  6.1\% &  66.0\% \\
\textbf{Bot}   & \textbf{1\,441} & \textbf{1.0\%} & \textbf{23.5\%} \\
\bottomrule
\end{tabular}
\end{table}

%-----------------------------------------------------------------------
\section{Failure Analysis}
%-----------------------------------------------------------------------

Model~C misses 19\,838 flows (5.90\% of attacks) with 23\,272 false
positives (6.13\% FPR).  All missed attacks share a common property:
both $s_{\text{GMM}} \approx 0.01$--$0.04$ and $s_{\text{LSTM}} \approx
0.15$--$0.45$---both scores are well below the decision boundary.

\textbf{DoS Hulk} (11\,627 FN, 93.3\% detected): Hulk generates
randomised HTTP requests; each flow's packet sizes and ports fall within
normal bounds.  GMM $\approx 0.024$ (marginally normal), LSTM $\approx
0.147$ (sequence looks like bursty web traffic).

\textbf{Bot} (1\,102 FN, 23.5\% detected): C2 beaconing with 60+\,s
inter-session intervals.  At network rates, 50 flows span $<$1\,s, so
the W=50 window never captures the beaconing rhythm.  GMM $\approx 0.026$,
LSTM $\approx 0.156$.  Detection requires W$\gg$50 or session-level
aggregation over hours.

\textbf{Fix path}: (i) application-layer payload features for HTTP-based
attacks; (ii) W$\,\gg\,$50 or graph-based features for Bot.

%-----------------------------------------------------------------------
\section{Concept Drift Detection}
%-----------------------------------------------------------------------

\subsection{Motivation}

A GMM trained on Monday's traffic may become stale by Friday.  If benign
traffic patterns change (new applications, firmware updates), the model's
false positive rate rises without warning.

\subsection{Page-Hinckley Test}

We monitor the running mean of GMM log-likelihood $\ell_t =
\log p_{\text{GMM}}(\mathbf{x}_t)$.  The Page-Hinckley statistic is:
\begin{align}
  \bar{\mu}_t &= \frac{1}{t}\sum_{i=1}^{t} \ell_i \\
  PH_t &= \max_{k\le t}\,\bar{\mu}_k - \bar{\mu}_t
\end{align}
Drift is declared when $PH_t > \lambda$.  After detection, the running
mean resets to the current value.

\begin{algorithm}[h]
\caption{Page-Hinckley Drift Detection}
\begin{algorithmic}[1]
\Require Baseline scores $\{b_i\}$, threshold $\lambda$, min\_instances $m$
\State $\bar{\mu} \gets \text{mean}(\{b_i\})$,\quad $\mu_{\max} \gets \bar{\mu}$,\quad $n \gets 0$
\For{each new score $\ell$}
  \State $n \gets n+1$;\quad $\bar{\mu} \gets \bar{\mu} + (\ell - \bar{\mu})/n$
  \State $\mu_{\max} \gets \max(\mu_{\max},\, \bar{\mu})$
  \State $PH \gets \mu_{\max} - \bar{\mu}$
  \If{$n > m$ \textbf{and} $PH > \lambda$}
    \State \textbf{emit} drift event;\quad reset $\bar{\mu} \gets \ell$,\; $\mu_{\max} \gets \ell$,\; $n \gets 0$
  \EndIf
\EndFor
\end{algorithmic}
\end{algorithm}

\textbf{Complexity}: O(1) time and O(1) space per flow---suitable for
line-rate stream processing.

\subsection{Results}

Fitted on 5\,000 benign training flows (baseline mean LL\,=\,$-28.6$).
Run on 716\,092 temporal test flows with $\lambda=50$.
\textbf{1 drift event detected at flow 379\,600}---13 flows into the
FTP-Patator attack period, immediately after the benign-to-attack
transition.

Sensitivity analysis:
\begin{table}[h]
\centering
\caption{PH Test Sensitivity ($\lambda$)}
\label{tab:ph}
\begin{tabular}{ccc}
\toprule
$\lambda$ & Drift events & First detection \\
\midrule
10 & 15 & flow 502 \\
20 &  3 & flow 22\,952 \\
30 &  2 & flow 35\,527 \\
50 &  1 & flow 379\,600 \\
100 & 0 & -- \\
\bottomrule
\end{tabular}
\end{table}

$\lambda=50$ provides a good trade-off: it ignores transient anomalies
(single-flow spikes) but detects the sustained distribution shift at the
benign$\to$attack boundary.

%-----------------------------------------------------------------------
\section{Conclusion}
%-----------------------------------------------------------------------

Model~C achieves F1=93.6\% and AUC=98.7\% on CICIDS-2017, outperforming
GMM alone on both metrics.  The F1 journey across all phases:
\begin{center}
$50.1\%$ (baseline) $\to$ $90.97\%$ (GMM) $\to$ $63.94\%$ (LSTM shuffled)
$\to$ $60.2\%$ (LSTM temporal) $\to$ $\mathbf{93.6\%}$ (Hybrid~C)
\end{center}

Remaining limitations: (i) Bot detection at only 23.5\%---requires
session-level temporal features; (ii) HTTP-level attacks (Hulk, Web
Attacks) require payload inspection beyond 34 flow statistics;
(iii) the concept drift detector would retrain on any sustained
anomaly---adversarial triggering of retraining is a known risk.

\textbf{Future work}: online LSTM-AE fine-tuning on confirmed benign
windows; graph-based flow aggregation for Bot/Infiltration; longer windows
(W=500) for slow beaconing.

\begin{thebibliography}{9}

\bibitem{cicids}
I.~Sharafaldin, A.~H.~Lashkari, and A.~A.~Ghorbani, ``Toward Generating a
New Intrusion Detection Dataset and Intrusion Traffic Characterization,''
\textit{ICISSP}, 2018.

\bibitem{hochreiter}
S.~Hochreiter and J.~Schmidhuber, ``Long Short-Term Memory,''
\textit{Neural Computation}, vol.~9, no.~8, 1997.

\bibitem{pagehinckley}
E.~S.~Page, ``Continuous Inspection Schemes,'' \textit{Biometrika},
vol.~41, no.~1/2, 1954.

\bibitem{bishop}
C.~M.~Bishop, \textit{Pattern Recognition and Machine Learning}.
Springer, 2006.

\bibitem{breiman}
L.~Breiman, ``Random Forests,'' \textit{Machine Learning},
vol.~45, no.~1, 2001.

\end{thebibliography}

\end{document}
